\section{Code--Data Identification and Data-Preserving Rewriting}
\label{sec:code_data}
\subsection{Problem and Representations}
For an input byte sequence $B=(b_0,\ldots,b_{n-1})$, identification
returns candidate data intervals $[s_j,e_j)$ in payload-relative
coordinates, with optional kinds such as ASCII, UTF-16LE, or another
data object. Acceptance for rewriting additionally depends on reference
evidence and implementation support. The original inline-data and
immediate-operand examples in Tables~\ref{tab:broken-payload}
and~\ref{tab:functional-payload} motivate this added layer.

The implementation distinguishes a model's reported intervals, regions
accepted by deterministic checks, and regions actually extracted.
These sets need not coincide. An unsupported reference can prevent
processing of a correctly identified object, while filtering can make
different proposals yield identical final binaries. Evaluation retains
the output stage associated with each observation.

\subsection{Symbolic and Heuristic Identification}
The non-LLM path combines CFGFast candidate blocks and inferred data
with bounded symbolic exploration. Observed instruction starts are
expanded into their decoded byte spans, and newly observed entries
can seed a further graph-recovery pass. This prevents interior bytes
of an observed instruction from being treated as unvisited data solely
because only its start address was recorded.

The initial state leaves uninitialized registers and memory symbolic;
it does not reproduce a complete Windows process or all API behavior.
Time, state/block-step, pending-state, and solver limits bound the
analysis. Reports record termination reasons, unresolved states, and
errors. Missing observations therefore do not prove that a region is
data or unreachable.

Printable ASCII and UTF-16LE patterns supply additional candidate
evidence. Observed code spans prevent their acceptance under the
string heuristic. For strings overlapping candidate graph code,
additional structural corroboration is required. Graph-inferred
data and bytes outside recovered blocks retain different evidential
categories; the latter are uncertain rather than a complete data
partition. The available analysis limits and candidate categories
accompany the classification result.

\subsection{LLM Identification and Deterministic Handling}
The alternative path presents hex bytes, a heuristic linear disassembly,
and printable candidates to the model. It requests structured intervals
instead of rewritten code. The parser checks bounds, positive lengths,
and overlap. Reference scanning and data handling remain deterministic;
a proposal is applied to a fresh engine instance so that rejection
does not reuse partially modified state.

Accepted data contents and supported references are retained through
layout changes. The engine finalizes the supported references once
the code layout is known and checks integrity of the preserved data.
These checks do not establish that every reference in an arbitrary
input has been recovered. Architecture-specific operands require
separate x64 and x86 implementations even though the conceptual
contract is shared.

An invalid proposal can trigger feedback and another request in the
general pipeline. The reported classifier experiments specify their
actual call budgets separately from the transformation experiments.
The classifier's default worked-example condition is not labeled
zero-shot or two-shot merely because those conditions are compared
in assembly transformation.

Caching reuses identified regions for an input identity while references
are recomputed for the working representation. Timing must distinguish
new analysis from cached results. On curated fixtures, independent
annotations allow reference-based assessment; on the selected real
payloads, agreement and functional checks remain different forms of
evidence. The experiments in Section~\ref{sec:extended_results}
explicitly retain that distinction.
